Python 中与列表和序列天然互补的，是字典（dictionary）类型，它将查找键映射到对应的值（通常称为关联数组（associative array）或哈希表（hash table））。字典的通用性，使其成为记账式管理的理想之选：动态地跟踪新增与不断变化的数据片段，以及它们之间的相互关系。在编写一个新程序时，在还不确定需要哪些其他数据结构或类之前，使用字典是一个很好的起点。

字典在添加与删除条目时提供常数时间（摊还）的性能，远胜于简单列表自身所能达到的水平。因此，字典作为 Python 用以实现其面向对象特性的核心数据结构，也就不难理解了。Python 还有特殊的语法与相关的内置模块，能为字典赋予超出你在其他语言中从简单的哈希表类型所能预期之外的额外能力。

\section{Item 25: 依赖字典插入顺序时需谨慎}

在 Python 3.5 及更早版本中，遍历一个字典（dictionary）实例会以其任意顺序返回键。遍历的顺序与条目最初插入字典的顺序并不一致。例如，这里我用 Python 3.5 版本创建了一个字典，把动物名称映射到它们对应的幼崽名称：

\begin{codebox}[python]
# Python 3.5
baby_names = {
    "cat": "kitten",
    "dog": "puppy",
}
print(baby_names)
>>>
{'dog': 'puppy', 'cat': 'kitten'}
\end{codebox}

当我创建这个字典时，键的顺序是 \texttt{"cat"}、\texttt{"dog"}，但当我打印它时，键的顺序却反了过来：\texttt{"dog"}、\texttt{"cat"}。这种行为令人意外，使得测试用例更难复现，增加了调试的难度，对 Python 新手尤其容易困惑。

这种情况之所以发生，是因为旧版 Python 中的字典（dictionary）类型实现其哈希表算法时，结合了内置（built-in）的 \texttt{hash} 函数，以及一个在 Python 解释器（interpreter）进程开始执行时分配的随机种子。这些行为共同作用，导致字典的顺序既与插入顺序不一致，又会在每次程序执行之间随机打乱。

从 Python 3.6 开始，字典开始保留插入顺序；而自 3.7 版本起，这正式成为 Python 语言规范的一部分。现在，下面这段代码总会以程序员当初创建字典时的相同顺序打印它：

\begin{codebox}[python]
baby_names = {
    "cat": "kitten",
    "dog": "puppy",
}
print(baby_names)
>>>
{'cat': 'kitten', 'dog': 'puppy'}
\end{codebox}

在 Python 3.5 及更早版本中，\texttt{dict} 类型提供的所有依赖遍历顺序的方法——包括 \texttt{keys}、\texttt{values}、\texttt{items} 和 \texttt{popitem}——同样表现出这种看似随机的行为：

\begin{codebox}[python]
# Python 3.5
print(list(baby_names.keys()))
print(list(baby_names.values()))
print(list(baby_names.items()))
print(baby_names.popitem())  # 随机选取一项
>>>
['dog', 'cat']
['puppy', 'kitten']
[('dog', 'puppy'), ('cat', 'kitten')]
('dog', 'puppy')
\end{codebox}

现在，这些方法都提供了一致的插入顺序，你在编写程序时可以放心依赖：

\begin{codebox}[python]
print(list(baby_names.keys()))
print(list(baby_names.values()))
print(list(baby_names.items()))
print(baby_names.popitem())  # 最后插入的项
>>>
['cat', 'dog']
['kitten', 'puppy']
[('cat', 'kitten'), ('dog', 'puppy')]
('dog', 'puppy')
\end{codebox}

这一变化对依赖 \texttt{dict} 类型及其具体实现的其他 Python 特性产生了许多影响。函数的关键字参数（keyword argument）——包括捕获所有参数的 \texttt{**kwargs} 参数（参见 Item 35（"用关键字参数提供可选行为"）与 Item 37（"用仅限关键字与仅限位置参数确保清晰"））——在更早的 Python 版本中会以看似随机的顺序传入，这会让函数调用的调试变得更加困难：

\begin{codebox}[python]
# Python 3.5
def my_func(**kwargs):
    for key, value in kwargs.items():
        print("%s = %s" % (key, value))
my_func(goose="gosling", kangaroo="joey")
>>>
kangaroo = joey
goose = gosling
\end{codebox}

而现在在最新版本的 Python 中，关键字参数的顺序总是被保留，以匹配程序员最初调用函数时的顺序：

\begin{codebox}[python]
def my_func(**kwargs):
    for key, value in kwargs.items():
        print(f"{key} = {value}")
my_func(goose="gosling", kangaroo="joey")
>>>
goose = gosling
kangaroo = joey
\end{codebox}

类（class）也使用 \texttt{dict} 类型来存储其实例字典（instance dictionary）。在更早的 Python 版本中，\texttt{object} 的字段会表现出这种随机化行为：

\begin{codebox}[python]
# Python 3.5
class MyClass:
    def __init__(self):
        self.alligator = "hatchling"
        self.elephant = "calf"
a = MyClass()
for key, value in a.__dict__.items():
    print("%s = %s" % (key, value))
>>>
elephant = calf
alligator = hatchling
\end{codebox}

同样，你现在可以放心假定这些实例字段的赋值顺序会反映在 \texttt{\_\_dict\_\_} 中：

\begin{codebox}[python]
class MyClass:
    def __init__(self):
        self.alligator = "hatchling"
        self.elephant = "calf"
a = MyClass()
for key, value in a.__dict__.items():
    print(f"{key} = {value}")
>>>
alligator = hatchling
elephant = calf
\end{codebox}

字典保留插入顺序的方式如今已成为 Python 语言规范的一部分。对于上述语言特性，你可以依赖这种行为，甚至可以把它作为你为类和函数设计的 API 的一部分（示例参见 Item 65（"考虑用类体定义顺序建立属性间关系"））。

注意：长期以来，内置的 \texttt{collections} 模块一直有一个 \texttt{OrderedDict} 类，它也保留插入顺序。尽管这个类的行为自 Python 3.7 起与标准 \texttt{dict} 类型相似，但 \texttt{OrderedDict} 的性能特征却大不相同。如果你需要处理大量的键插入和 \texttt{popitem} 调用（例如，实现一个最近最少使用（least-recently-used, LRU）缓存），\texttt{OrderedDict} 可能比标准的 Python \texttt{dict} 类型更合适（关于如何确认是否真的需要它，参见 Item 92（"在优化之前先做性能分析"））。

不过，在处理字典时，你不应该总是假定插入顺序行为一定存在。Python 让程序员很容易定义自己的自定义容器类型，去模拟匹配 \texttt{list}、\texttt{dict} 及其他类型的标准协议（参见 Item 57（"从 collections.abc 类继承以定义自定义容器类型"））。Python 不是静态类型（statically typed）的语言，因此大多数代码依赖鸭子类型（duck typing）——即对象的行为就是它的事实类型——而不是僵化的类层次结构（参见 Item 3（"永远别指望 Python 在编译期检测错误"））。这可能会导致一些令人意外的陷阱。

举例来说，假设我在编写一个程序来展示"最可爱幼崽动物"比赛的结果。这里，我先从一个字典开始，它包含每只动物获得的总票数：

\begin{codebox}[python]
votes = {
    "otter": 1281,
    "polar bear": 587,
    "fox": 863,
}
\end{codebox}

现在，我定义一个函数来处理这份投票数据，并把每只动物名称的排名保存到传入的一个空字典中。在这种情况下，这个字典可能是驱动某个 UI 元素的数据模型：

\begin{codebox}[python]
def populate_ranks(votes, ranks):
    names = list(votes.keys())
    names.sort(key=votes.get, reverse=True)
    for i, name in enumerate(names, 1):
        ranks[name] = i
\end{codebox}

我还需要一个函数，告诉我哪只动物赢得了比赛。这个函数之所以能工作，是因为它假定 \texttt{populate\_ranks} 会按升序为 \texttt{ranks} 字典赋值，也就是说第一个键必定是获胜者：

\begin{codebox}[python]
def get_winner(ranks):
    return next(iter(ranks))
\end{codebox}

在这里，我确认这些函数按设计工作，并给出我所期望的结果：

\begin{codebox}[python]
ranks = {}
populate_ranks(votes, ranks)
print(ranks)
winner = get_winner(ranks)
print(winner)
>>>
{'otter': 1, 'fox': 2, 'polar bear': 3}
otter
\end{codebox}

现在，想象这个程序的需求发生了变化。显示结果的 UI 元素应当按字母顺序而非排名顺序排列。为此，我可以使用内置的 \texttt{collections.abc} 模块来定义一个新的、类似字典的类，它按字母顺序遍历其内容：

\begin{codebox}[python]
from collections.abc import MutableMapping
class SortedDict(MutableMapping):
    def __init__(self):
        self.data = {}
    def __getitem__(self, key):
        return self.data[key]
    def __setitem__(self, key, value):
        self.data[key] = value
    def __delitem__(self, key):
        del self.data[key]
    def __iter__(self):
        keys = list(self.data.keys())
        keys.sort()
        for key in keys:
            yield key
    def __len__(self):
        return len(self.data)
\end{codebox}

我可以把一个 \texttt{SortedDict} 实例替换标准的 \texttt{dict} 来使用之前那些函数，而且不会抛出任何错误，因为这个类遵循了标准字典的协议。然而，结果却是错误的：

\begin{codebox}[python]
sorted_ranks = SortedDict()
populate_ranks(votes, sorted_ranks)
print(sorted_ranks.data)
winner = get_winner(sorted_ranks)
print(winner)
>>>
{'otter': 1, 'fox': 2, 'polar bear': 3}
fox
\end{codebox}

这里的问题在于，\texttt{get\_winner} 的实现假定字典的遍历是按插入顺序进行的，以匹配 \texttt{populate\_ranks}。但这段代码用的是 \texttt{SortedDict} 而非 \texttt{dict}，所以这个假定不再成立。于是，返回给获胜者的值是 \texttt{"fox"}，也就是字母顺序排在最前面的那个。

有三种方法可以缓解这一问题。首先，我可以重新实现 \texttt{get\_winner} 函数，使其不再假定 \texttt{ranks} 字典具有某种特定的遍历顺序。这是最保守、最稳健的方案：

\begin{codebox}[python]
def get_winner(ranks):
    for name, rank in ranks.items():
        if rank == 1:
            return name
winner = get_winner(sorted_ranks)
print(winner)
>>>
otter
\end{codebox}

第二种做法是在函数开头添加一个显式检查，以确保 \texttt{ranks} 的类型符合我的预期；如果不符合，就抛出一个异常（exception）。这种方案的运行（runtime）性能很可能优于更保守的做法：

\begin{codebox}[python]
def get_winner(ranks):
    if not isinstance(ranks, dict):
        raise TypeError("must provide a dict instance")
    return next(iter(ranks))
get_winner(sorted_ranks)
>>>
Traceback ...
TypeError: must provide a dict instance
\end{codebox}

第三种替代方案是使用类型标注（type annotation）来强制要求传给 \texttt{get\_winner} 的值是一个 \texttt{dict} 实例，而不是一个具有字典行为的 \texttt{MutableMapping}（参见 Item 124（"考虑通过 typing 做静态分析以消除缺陷"））。在这里，我在上述带类型标注的代码版本上用严格模式运行了 \texttt{mypy} 工具：

\begin{codebox}[python]
from typing import Dict, MutableMapping
def populate_ranks( votes: Dict[str, int],
ranks: Dict[str, int]) -> None:
    names = list(votes.keys())
    names.sort(key=votes.__getitem__, reverse=True)
    for i, name in enumerate(names, 1):
        ranks[name] = i
def get_winner(ranks: Dict[str, int]) -> str:
    return next(iter(ranks))
class SortedDict(MutableMapping[str, int]):
    ...
votes = {
    "otter": 1281,
    "polar bear": 587,
    "fox": 863,
}
sorted_ranks = SortedDict()
populate_ranks(votes, sorted_ranks)
print(sorted_ranks.data)
winner = get_winner(sorted_ranks)
print(winner)
$ python3 -m mypy --strict example.py
.../example.py:48: error: Argument 2 to "populate_ranks" has
incompatible type "SortedDict"; expected "dict[str, int]"
[arg-type]
.../example.py:50: error: Argument 1 to "get_winner" has
incompatible type "SortedDict"; expected "dict[str, int]"
[arg-type]
Found 2 errors in 1 file (checked 1 source file)
\end{codebox}

这正确地检测出了 \texttt{dict} 与 \texttt{SortedDict} 类型之间的不匹配，并把不正确的用法标记为错误。这种方案在静态类型安全与运行（runtime）性能之间提供了最佳平衡。

\subsection{需要记住的要点}

\begin{itemize}
\item 自 Python 3.7 起，你可以放心依赖这样一个事实：遍历一个字典实例的内容时，顺序与它最初添加键的顺序相同。
\item Python 让定义行为像字典、却又不是 \texttt{dict} 实例的对象变得很容易。对于这些类型，你不能假定插入顺序会被保留。
\item 对于类似字典的类，有三种谨慎应对的方法：编写不依赖插入顺序的代码、在运行时显式检查是否为 \texttt{dict} 类型，或者使用类型标注与静态分析（static analysis）要求传入 \texttt{dict} 值。
\end{itemize}

\section{Item 26: 处理缺失的字典键时，优先用 get 而非 in 与 KeyError}

与字典（dictionary）交互有三项基本操作：访问、赋值和删除键（key）及其关联的值（value）。字典的内容是动态的，因此当你试图访问或删除某个键时，它完全可能——甚至很可能——并不存在。

例如，假设我想确定人们最喜欢的面包种类，以便为一家三明治店设计菜单。这里，我定义了一个由计数器（counter）组成的字典，记录每种面包当前获得的票数：

\begin{codebox}[python]
counters = {
    "pumpernickel": 2,
    "sourdough": 1,
}
\end{codebox}

要为一次新投票增加计数，我需要先判断该键是否存在；若不存在，就以默认的零值插入该键；然后再将计数器的值加一。这需要访问键两次、赋值一次。下面我通过使用一个 \texttt{if} 语句配合 \texttt{in} 运算符来完成这项任务，当键存在时，\texttt{in} 运算符会返回 \texttt{True}：

\begin{codebox}[python]
key = "wheat"
if key in counters:
    count = counters[key]
else:
    count = 0
counters[key] = count + 1
print(counters)
>>>
{'pumpernickel': 2, 'sourdough': 1, 'wheat': 1}
\end{codebox}

实现相同行为的另一种方式是利用字典在尝试获取不存在键的值时会抛出 \texttt{KeyError} 异常这一点。暂且不论抛出异常与捕获异常的开销（参见 Item 80（「充分利用 \texttt{try}/\texttt{except}/\texttt{else}/\texttt{finally} 中的每一个代码块」）），这种方法在理论上更高效，因为它只需要一次访问和一次赋值：

\begin{codebox}[python]
try:
    count = counters[key]
except KeyError:
    count = 0
counters[key] = count + 1
\end{codebox}

这种「取出已存在的键，或返回默认值」的流程非常常见，以至于内置（built-in）类型 \texttt{dict} 专门提供了 \texttt{get} 方法来完成这一任务。\texttt{get} 的第二个参数是在键（第一个参数）不存在时要返回的默认值（关于这是否是一个良好的接口，参见 Item 32（「优先抛出异常而非返回 \texttt{None}」））。这种方法同样只需要一次访问和一次赋值，但它比 \texttt{KeyError} 的例子短得多，并且避免了异常处理的开销：

\begin{codebox}[python]
count = counters.get(key, 0)
counters[key] = count + 1
\end{codebox}

可以通过各种方式缩短 \texttt{in} 运算符和 \texttt{KeyError} 这两种写法，但所有这些替代方案都存在一个缺点：赋值语句需要重复书写，这降低了可读性，因此不值得采用：

\begin{codebox}[python]
if key not in counters:
    counters[key] = 0
counters[key] += 1
if key in counters:
    counters[key] += 1
else:
    counters[key] = 1
try:
    counters[key] += 1
except KeyError:
    counters[key] = 1
\end{codebox}

因此，对于存储简单类型的字典，使用 \texttt{get} 方法是最简短、最清晰的选择。

注意：如果你正在维护类似这样的计数器字典，不妨考虑内置模块 \texttt{collections} 中的 \texttt{Counter} 类，它提供了你大概需要的大部分功能。

如果字典中的值是像列表（list）这样更复杂的类型呢？例如，假设除了统计票数，我还想知道每种面包分别是谁投的票。这里，我通过将一个姓名列表与每个键关联起来实现这一点：

\begin{codebox}[python]
votes = {
    "baguette": ["Bob", "Alice"],
    "ciabatta": ["Coco", "Deb"],
}
key = "brioche"
who = "Elmer"
if key in votes:
    names = votes[key]
else:
    votes[key] = names = []
names.append(who)
print(votes)
>>>
{'baguette': ['Bob', 'Alice'],
 'ciabatta': ['Coco', 'Deb'],
 'brioche': ['Elmer']}
\end{codebox}

依赖 \texttt{in} 运算符，在键存在时需要访问两次，在键缺失时需要访问一次并赋值一次。这个例子与上面的计数器例子不同，因为对于每个键，如果它尚不存在，其值可以直接盲目地赋为默认空列表（empty list）。这个三元赋值语句（\texttt{votes[key] = names = []}）在一行内就补全了该键，而不需要两行。一旦默认值被插入字典，我就无需再次赋值，因为后续的 \texttt{append} 调用会通过引用（reference）修改该列表（背景知识参见 Item 30（「了解函数参数可能被修改」））。

当字典的值为列表时，也可以依赖抛出 \texttt{KeyError} 异常。这种方法在键存在时只需访问键一次，在键缺失时只需访问键一次并赋值一次，因此比 \texttt{in} 运算符更高效（忽略异常处理机制的开销）：

\begin{codebox}[python]
try:
    names = votes[key]
except KeyError:
    votes[key] = names = []
names.append(who)
\end{codebox}

类似地，我可以使用 \texttt{get} 方法在键存在时取出列表值，或者在键不存在时进行一次取值和一次赋值：

\begin{codebox}[python]
names = votes.get(key)
if names is None:
    votes[key] = names = []
names.append(who)
\end{codebox}

如果在 \texttt{if} 语句中使用赋值表达式（assignment expression，Python 3.8 引入；参见 Item 8（「用赋值表达式避免重复」）），那么这种用 \texttt{get} 取出列表值的方法还能再缩短一行，从而提升可读性：

\begin{codebox}[python]
if (names := votes.get(key)) is None:
    votes[key] = names = []
names.append(who)
\end{codebox}

值得注意的是，\texttt{dict} 类型还提供了 \texttt{setdefault} 方法，能进一步缩短这种模式。\texttt{setdefault} 尝试取出字典中某个键的值。如果该键不存在，该方法会将该键赋为所提供的默认值。然后，该方法返回该键对应的值：要么是原本就存在的值，要么是新插入的默认值。下面，我使用 \texttt{setdefault} 实现与上面 \texttt{get} 例子相同的逻辑：

\begin{codebox}[python]
names = votes.setdefault(key, [])
names.append(who)
\end{codebox}

这按预期工作，而且比使用带赋值表达式的 \texttt{get} 更短。然而，这种方法的可读性并不理想。方法名 \texttt{setdefault} 并没有让它的用途一目了然。明明是在「取」一个值，为什么叫 \texttt{set}？为什么不叫 \texttt{get\_or\_set}？我在这里有点吹毛求疵（arguing about the color of the bike shed），但重点是：如果你是一个刚接触这段代码、对 Python 并不完全熟悉的新读者，你可能会因为 \texttt{setdefault} 不够自解释而难以理解这段代码想要做什么。

还有一个重要的陷阱（gotcha）：传给 \texttt{setdefault} 的默认值在键缺失时会被直接赋入字典，而不是被复制。下面，我以值为列表的情形演示这一影响：

\begin{codebox}[python]
data = {}
key = "foo"
value = []
data.setdefault(key, value)
print("Before:", data)
value.append("hello")
print("After: ", data)
>>>
Before: {'foo': []}
After:  {'foo': ['hello']}
\end{codebox}

因此，我必须确保每次用 \texttt{setdefault} 访问键时，都构造一个新的默认值。这在这个例子中会带来显著的性能（performance）开销，因为我必须为每次调用都分配一个 \texttt{list} 实例。如果我复用同一个对象作为默认值——为了提升效率或可读性，我可能会这样做——就可能引入奇怪的行为和缺陷（bug）（此问题的另一个例子参见 Item 36（「用 \texttt{None} 和文档字符串指定动态默认参数」））。

回到之前那个用计数器而非投票人列表作为字典值的例子：为什么不在那种情况下也使用 \texttt{setdefault} 方法呢？下面，我用这种方法重新实现同一例子：

\begin{codebox}[python]
count = counters.setdefault(key, 0)
counters[key] = count + 1
\end{codebox}

这里的问题在于，调用 \texttt{setdefault} 是多余的。在计数器加一之后，你总是需要把字典中的键赋为一个新值，因此 \texttt{setdefault} 所做的额外赋值是没必要的。前面那种用 \texttt{get} 更新计数器的方法只需要一次访问和一次赋值，而使用 \texttt{setdefault} 则需要一次访问和两次赋值。

只有在极少数情况下，使用 \texttt{setdefault} 才是处理缺失字典键最短的方式（例如，构造代价很低且不会抛出异常 \texttt{list} 实例默认值）。在这些非常特殊的情形下，为了避免多写字符和行数而使用 \texttt{get}，似乎值得接受这个令人困惑的方法名 \texttt{setdefault}。然而，在这些情况下你真正应该做的，往往是改用 \texttt{defaultdict}（参见 Item 27（「处理内部状态中缺失的项时，优先用 \texttt{defaultdict} 而非 \texttt{setdefault}」））。

\subsection{需要记住的要点}

\begin{itemize}
\item 检测并处理字典中缺失的键有四种常见方式：使用 \texttt{in} 运算符、\texttt{KeyError} 异常、\texttt{get} 方法，以及 \texttt{setdefault} 方法。
\item \texttt{get} 方法最适合包含计数器这类基本类型的字典；而当构造字典默认值代价很高或可能抛出异常时，配合赋值表达式使用 \texttt{get} 方法更为可取。
\item 当 \texttt{dict} 的 \texttt{setdefault} 方法看起来最适合你的问题时，你应当转而考虑使用 \texttt{defaultdict}。
\end{itemize}

\section{Item 27: 处理内部状态中缺失的项时，优先用 defaultdict 而非 setdefault}

当你处理一个并非由你创建的字典（dictionary）时，有多种方式可以处理缺失的键（参见 Item 26（「处理缺失的字典键时，优先用 \texttt{get} 而非 \texttt{in} 与 \texttt{KeyError}」））。尽管使用 \texttt{get} 方法比使用 \texttt{in} 运算符和 \texttt{KeyError} 异常更好，但在某些使用场景下，\texttt{setdefault} 似乎是最短的选择。

例如，假设我想记录自己到访过世界各国的哪些城市。这里，我使用一个字典来实现：它将国家名映射到包含对应城市名的 \texttt{set} 实例：

\begin{codebox}[python]
visits = {
    "Mexico": {"Tulum", "Puerto Vallarta"},
    "Japan": {"Hakone"},
}
\end{codebox}

无论国家名是否已经存在于字典中，我都可以使用 \texttt{setdefault} 方法向这些集合中添新的城市。这段代码比用 \texttt{get} 方法配合赋值表达式实现相同行为要短得多：

\begin{codebox}[python]
# Short
visits.setdefault("France", set()).add("Arles")
# Long
if (japan := visits.get("Japan")) is None:
    visits["Japan"] = japan = set()
japan.add("Kyoto")
print(visits)
>>>
{'Mexico': {'Puerto Vallarta', 'Tulum'},
 'Japan': {'Kyoto', 'Hakone'},
 'France': {'Arles'}}
\end{codebox}

那么，当你确实掌控着被访问字典的创建时，情况又如何呢？通常，当你用一个字典实例来记录某个对象（object）的内部状态（internal state）时，就属于这种情况。下面，我用一个包含辅助方法的类把上面的例子包起来，以访问其存储在字典中的动态内部状态：

\begin{codebox}[python]
class Visits:
    def __init__(self):
        self.data = {}
    def add(self, country, city):
        city_set = self.data.setdefault(country, set())
        city_set.add(city)
\end{codebox}

这个新类隐藏了使用正确参数调用 \texttt{setdefault} 的复杂性，并为程序员提供了更友好的接口（interface）：

\begin{codebox}[python]
visits = Visits()
visits.add("Russia", "Yekaterinburg")
visits.add("Tanzania", "Zanzibar")
print(visits.data)
>>>
{'Russia': {'Yekaterinburg'}, 'Tanzania': {'Zanzibar'}}
\end{codebox}

然而，\texttt{Visits.add} 方法的实现并不理想。\texttt{setdefault} 方法的命名依然令人困惑，这使得代码的新读者更难立刻理解发生了什么。而且这种实现效率不高，因为它在每次调用时都会构造一个新的 \texttt{set} 实例，无论给定的国家是否已经存在于 \texttt{data} 字典中。

幸运的是，内置模块 \texttt{collections} 中的 \texttt{defaultdict} 类简化了这一常见用例：当键不存在时，它会自动存储一个默认值。你所要做的只是提供一个函数（function），在每次键缺失时返回要使用的默认值（这是 Item 48（「对于简单接口，优先接受函数而非类」）的一个例子）。下面，我重写 \texttt{Visits} 类以使用 \texttt{defaultdict}：

\begin{codebox}[python]
from collections import defaultdict
class Visits:
    def __init__(self):
        self.data = defaultdict(set)
    def add(self, country, city):
        self.data[country].add(city)
visits = Visits()
visits.add("England", "Bath")
visits.add("England", "London")
print(visits.data)
>>>
defaultdict(<class 'set'>, {'England': {'Bath', 'London'}})
\end{codebox}

现在，\texttt{add} 的实现简短而清晰。代码可以假定：访问 \texttt{data} 字典中的任何键，都必定会得到一个已存在的 \texttt{set} 实例。不会分配多余的 \texttt{set} 实例，而当 \texttt{add} 方法被调用很多次时，这种多余分配可能会代价高昂。

对于这类情况，使用 \texttt{defaultdict} 远比使用 \texttt{setdefault} 更好（另一个例子参见 Item 29（「用组合类代替深度嵌套的字典、列表与元组」））。仍有一些情形是 \texttt{defaultdict} 无法解决的，但 Python 中还有更多工具可以绕开这些限制（参见 Item 28（「了解如何用 \texttt{\_\_missing\_\_} 构造依赖于键的默认值」）、Item 57（「为自定义容器类型继承自 \texttt{collections.abc} 类」），以及内置类 \texttt{collections.Counter}）。

\subsection{需要记住的要点}

\begin{itemize}
\item 如果你正在创建一个字典来管理任意一组潜在的键，那么只要 \texttt{defaultdict} 适合你的问题，就应当优先使用内置模块 \texttt{collections} 中的 \texttt{defaultdict} 实例。
\item 如果别人传给你一个包含任意键的字典，而你并不掌控它的创建，那么访问其元素时应当优先使用 \texttt{get} 方法。不过，在少数能因此写出更短代码、且默认对象分配开销很低的情形下，也值得考虑使用 \texttt{setdefault} 方法。
\end{itemize}

\section{Item 28: 用 \texttt{\_\_missing\_\_} 构造依赖键的默认值}

内置（built-in）\texttt{dict} 类型的 \texttt{setdefault} 方法，在特定情况下处理缺失键（missing key）时能写出更简短的代码（参见 Item 26（「用 \texttt{get} 替代 \texttt{in} 和 \texttt{KeyError} 来处理缺失的字典键」））。但在许多这类场景中，更合适的工具其实是内置模块 \texttt{collections} 中的 \texttt{defaultdict} 类型（原因见 Item 27（「用 \texttt{defaultdict} 替代 \texttt{setdefault} 来处理内部状态里缺失的项」））。不过，有些时候 \texttt{setdefault} 和 \texttt{defaultdict} 都不太合适。

举例来说，假设我要写一个程序来管理文件系统上的社交网络头像。我需要一个字典，把头像文件路径名映射到已打开的文件句柄（file handle），这样就能按需读写这些图片。在这里，我使用一个普通的 \texttt{dict} 实例，并用 \texttt{get} 方法和赋值表达式（assignment expression，即海象运算符 \texttt{:=}）来检查键是否存在（参见 Item 8（「用赋值表达式避免重复」））：

\begin{codebox}[python]
pictures = {}
path = "profile_1234.png"
if (handle := pictures.get(path)) is None:
    try:
        handle = open(path, "a+b")
    except OSError:
        print(f"Failed to open path {path}")
        raise
    else:
        pictures[path] = handle
handle.seek(0)
image_data = handle.read()
\end{codebox}

当文件句柄已经存在于字典中时，这段代码只做了一次字典访问。如果文件句柄不存在，则先用 \texttt{get} 访问一次字典，然后在 \texttt{try}/\texttt{except} 语句的 \texttt{else} 子句里赋值（参见 Item 80（「充分利用 \texttt{try}/\texttt{except}/\texttt{else}/\texttt{finally} 中的各个代码块」））。对 \texttt{read} 方法的调用，与调用 \texttt{open} 和处理异常的代码清楚地分离开来。

虽然也可以用 \texttt{in} 运算符或 \texttt{KeyError} 的方式来实现同样的逻辑，但这些做法需要更多的字典访问和更深层的嵌套。鉴于其他几种方式都能用，你大概会以为 \texttt{setdefault} 方法也能奏效：

\begin{codebox}[python]
try:
    handle = pictures.setdefault(path, open(path, "a+b"))
except OSError:
    print(f"Failed to open path {path}")
    raise
else:
    handle.seek(0)
    image_data = handle.read()
\end{codebox}

这段代码问题很多。用于创建文件句柄的内置函数 \texttt{open} 总是会被调用，即便 \texttt{path} 已经存在于字典中。这会导致多出一个文件句柄，可能与同一程序里已有的打开句柄相冲突。\texttt{open} 调用可能抛出的异常需要处理，但可能无法把它和同一行里 \texttt{setdefault} 调用可能抛出的异常区分开（对于其他类似字典的实现，这种情况是可能发生的；参见 Item 57（「通过继承 \texttt{collections.abc} 的基类来实现自定义容器类型」））。

如果你要管理内部状态，另一个可能产生的假设是：可以用 \texttt{defaultdict} 来追踪这些头像。下面我尝试用同样的思路来实现，但改用辅助函数（helper function）和 \texttt{defaultdict} 类：

\begin{codebox}[python]
from collections import defaultdict
def open_picture(profile_path):
    try:
        return open(profile_path, "a+b")
    except OSError:
        print(f"Failed to open path {profile_path}")
        raise
pictures = defaultdict(open_picture)
handle = pictures[path]
handle.seek(0)
image_data = handle.read()
>>>
Traceback ...
TypeError: open_picture() missing 1 required positional
argument: 'profile_path'
\end{codebox}

问题在于，\texttt{defaultdict} 要求传给其构造函数的函数不接受任何参数。这意味着 \texttt{defaultdict} 调用的辅助函数并不知道正在被访问的具体键（key），从而妨碍了我去调用 \texttt{open}。在这种情况下，\texttt{setdefault} 和 \texttt{defaultdict} 都无法满足我的需求。

好在这种情况足够常见，Python 还有另一种内置解法。你可以继承（subclass）\texttt{dict} 类型，并实现 \texttt{\_\_missing\_\_} 特殊方法（special method），为处理缺失键添加自定义逻辑。下面我通过定义一个新类来做到这一点，它复用了上面定义的同一个 \texttt{open\_picture} 辅助方法：

\begin{codebox}[python]
class Pictures(dict):
    def __missing__(self, key):
        value = open_picture(key)
        self[key] = value
        return value
pictures = Pictures()
handle = pictures[path]
handle.seek(0)
image_data = handle.read()
\end{codebox}

当 \texttt{pictures[path]} 这次字典访问发现 \texttt{path} 键不在字典中时，就会调用 \texttt{\_\_missing\_\_} 方法。该方法必须为这个键创建新的默认值，把它插入字典，并返回给调用者。之后再次访问同一个 \texttt{path} 就不会再调用 \texttt{\_\_missing\_\_} 了，因为对应的项已经存在（这与 \texttt{\_\_getattr\_\_} 的行为类似；参见 Item 61（「用 \texttt{\_\_getattr\_\_}、\texttt{\_\_getattribute\_\_} 和 \texttt{\_\_setattr\_\_} 实现惰性属性」））。

\subsection{需要记住的要点}
\begin{itemize}
\item 当创建默认值的计算代价很高，或者可能抛出异常时，\texttt{dict} 的 \texttt{setdefault} 方法并不合适。
\item 传给 \texttt{defaultdict} 的函数不能接受任何参数，这使得默认值无法依赖正在被访问的键。
\item 你可以定义自己的 \texttt{dict} 子类，并实现 \texttt{\_\_missing\_\_} 方法，从而构造那些必须知道访问了哪个键的默认值。
\end{itemize}

\section{Item 29: 用类组合代替深度嵌套的字典、列表与元组}

Python 的内置（built-in）字典（dictionary）类型非常适合在对象的整个生命周期内维护动态的内部状态。所谓动态，我指的是你需要为一组事先无法预料的标识符做记账（bookkeeping）的情形。例如，假设我想记录一组学生的成绩，而这些学生的姓名事先并不知道。我可以定义一个类，把姓名存进一个字典（dictionary），而不是给每个学生各用一个预定义的属性：

\begin{codebox}[python]
class SimpleGradebook:
    def __init__(self):
        self._grades = {}
    def add_student(self, name):
        self._grades[name] = []
    def report_grade(self, name, score):
        self._grades[name].append(score)
    def average_grade(self, name):
        grades = self._grades[name]
        return sum(grades) / len(grades)
\end{codebox}

使用这个类很简单：

\begin{codebox}[python]
book = SimpleGradebook()
book.add_student("Isaac Newton")
book.report_grade("Isaac Newton", 90)
book.report_grade("Isaac Newton", 95)
book.report_grade("Isaac Newton", 85)
print(book.average_grade("Isaac Newton"))
>>>
90.0
\end{codebox}

字典（dictionary）、列表（list）、元组（tuple）和集合（set）用起来太顺手了，以至于存在一种危险：它们会诱使你写出脆弱的代码。例如，假设我想扩展 \texttt{SimpleGradebook} 类，使其按科目（而不仅仅是总体）保存成绩的 \texttt{list}。为此，我可以把 \texttt{\_grades} 字典改为：把学生姓名（它的键）映射到另一个字典（它的值）。最内层的字典再把科目（它的键）映射到成绩列表（它的值）。这里我通过在内层字典上使用一个 \texttt{defaultdict} 实例来处理缺失的科目（背景知识参见 Item 27（"用 defaultdict 优于 setdefault 处理内部状态的缺失项"））：

\begin{codebox}[python]
from collections import defaultdict
class BySubjectGradebook:
    def __init__(self):
        self._grades = {}                       # Outer dict
    def add_student(self, name):
        self._grades[name] = defaultdict(list)  # Inner dict
    def report_grade(self, name, subject, grade):
        by_subject = self._grades[name]
        grade_list = by_subject[subject]
        grade_list.append(grade)
    def average_grade(self, name):
        by_subject = self._grades[name]
        total, count = 0, 0
        for grades in by_subject.values():
            total += sum(grades)
            count += len(grades)
        return total / count
\end{codebox}

这已经足够直观了。\texttt{report\_grade} 和 \texttt{average\_grade} 这两个方法为了处理多级字典而增加了不少复杂度，但看起来还算可控。使用这个类依然很简单：

\begin{codebox}[python]
book = BySubjectGradebook()
book.add_student("Albert Einstein")
book.report_grade("Albert Einstein", "Math", 75)
book.report_grade("Albert Einstein", "Math", 65)
book.report_grade("Albert Einstein", "Gym", 90)
book.report_grade("Albert Einstein", "Gym", 95)
print(book.average_grade("Albert Einstein"))
>>>
81.25
\end{codebox}

现在，假设需求再次发生变化。我还想跟踪每次成绩在班级总评中所占的权重，使得期中和期末考试比随堂小测更重要。实现这一功能的一种办法是修改最内层的字典：不再把科目（它的键）映射到成绩列表（它的值），而是在每个键对应的价值列表中使用 \texttt{(score, weight)} 这样的元组（tuple）。尽管对 \texttt{report\_grade} 的改动看起来很简单——只需让 \texttt{grade\_list} 存储元组（tuple）实例——但 \texttt{average\_grade} 方法现在变成了循环套循环，难以阅读：

\begin{codebox}[python]
class WeightedGradebook:
    def __init__(self):
        self._grades = {}
    def add_student(self, name):
        self._grades[name] = defaultdict(list)
    def report_grade(self, name, subject, score, weight):
        by_subject = self._grades[name]
        grade_list = by_subject[subject]
        grade_list.append((score, weight))    # Changed
    def average_grade(self, name):
        by_subject = self._grades[name]
        score_sum, score_count = 0, 0
        for scores in by_subject.values():
            subject_avg, total_weight = 0, 0
            for score, weight in scores:      # Added inner loop
                subject_avg += score * weight
                total_weight += weight
            score_sum += subject_avg / total_weight
            score_count += 1
        return score_sum / score_count
\end{codebox}

使用这个类也变得更困难了。那些位置参数（positional argument）中的数字究竟代表什么都不再清晰：

\begin{codebox}[python]
book = WeightedGradebook()
book.add_student("Albert Einstein")
book.report_grade("Albert Einstein", "Math", 75, 0.05)
book.report_grade("Albert Einstein", "Math", 65, 0.15)
book.report_grade("Albert Einstein", "Math", 70, 0.80)
book.report_grade("Albert Einstein", "Gym", 100, 0.40)
book.report_grade("Albert Einstein", "Gym", 85, 0.60)
print(book.average_grade("Albert Einstein"))
>>>
80.25
\end{codebox}

当你看到这样的复杂度时，就该从字典、列表、元组、集合这类内置（built-in）类型，跨越到由类构成的层次结构了。

在成绩这个例子里，起初我并不知道需要支持加权成绩，所以创建其他类的复杂度看起来并不划算。Python 的内置字典和元组类型让人很容易继续下去，一层又一层地往内部记账结构里添加。但你应当避免嵌套超过一层；使用包含字典的字典，会让其他程序员难以阅读你的代码，并为日后的维护埋下噩梦般的隐患（另一种处理方式参见 Item 9（"在流程控制中用 match 解构，当 if 语句已足够时避免使用"））。

一旦你意识到自己的记账逻辑变得复杂，就应当把它全部拆分成类。这样一来，你就能提供定义良好的接口（interface），更好地封装你的数据。这种做法还能让你在接口与具体实现之间创建出一层抽象（abstraction）。

\subsection{重构为类}

重构（refactoring）有多种做法（举例参见 Item 123（"考虑用 warnings 重构并迁移用法"））。在这个例子中，我可以从依赖树（dependency tree）的最底层开始改用类：单个成绩。对于这么简单的信息，用一个类显得过于重量级。不过，元组（tuple）倒是合适，因为成绩是不可变（immutable）的。这里我用 \texttt{(score, weight)} 形式的元组在列表中跟踪成绩：

\begin{codebox}[python]
grades = []
grades.append((95, 0.45))
grades.append((85, 0.55))
total = sum(score * weight for score, weight in grades)
total_weight = sum(weight for _, weight in grades)
average_grade = total / total_weight
\end{codebox}

我使用了 \texttt{\_}（下划线变量名，是 Python 中表示未使用变量的约定）来捕获每个成绩元组的第一个元素，并在计算 \texttt{total\_weight} 时忽略它。

这段代码的问题在于，元组（tuple）实例是依赖位置的。例如，如果我想给成绩关联更多信息（比如老师写的一组评语），就需要重写所有用到这个二元组的地方，让它们知道现在出现了三项而不是两项，这意味着我还得进一步使用 \texttt{\_} 来忽略某些下标（index）：

\begin{codebox}[python]
grades = []
grades.append((95, 0.45, "Great job"))
grades.append((85, 0.55, "Better next time"))
total = sum(score * weight for score, weight, _ in grades)
total_weight = sum(weight for _, weight, _ in grades)
average_grade = total / total_weight
\end{codebox}

这种把元组越拉越长的模式，与不断加深字典层级的做法如出一辙。一旦你发现自己用到了超过二元的元组，就该考虑换一种方式了。\texttt{dataclasses} 这个内置（built-in）模块恰好能满足我此时的需求：它让我能轻松地定义一个小巧的不可变类，把值存放在属性（attribute）中（参见 Item 56（"优先用 dataclasses 创建不可变对象"））：

\begin{codebox}[python]
from dataclasses import dataclass
@dataclass(frozen=True)
class Grade:
    score: int
    weight: float
\end{codebox}

接下来，我可以写一个类来表示单个科目，其中包含一组 \texttt{Grade} 实例：

\begin{codebox}[python]
class Subject:
    def __init__(self):
        self._grades = []
    def report_grade(self, score, weight):
        self._grades.append(Grade(score, weight))
    def average_grade(self):
        total, total_weight = 0, 0
        for grade in self._grades:
            total += grade.score * grade.weight
            total_weight += grade.weight
        return total / total_weight
\end{codebox}

然后，我写一个类来保存单个学生正在学习的科目集合：

\begin{codebox}[python]
class Student:
    def __init__(self):
        self._subjects = defaultdict(Subject)
    def get_subject(self, name):
        return self._subjects[name]
    def average_grade(self):
        total, count = 0, 0
        for subject in self._subjects.values():
            total += subject.average_grade()
            count += 1
        return total / count
\end{codebox}

最后，我写一个容器来容纳所有学生，并以姓名为键动态索引：

\begin{codebox}[python]
class Gradebook:
    def __init__(self):
        self._students = defaultdict(Student)
    def get_student(self, name):
        return self._students[name]
\end{codebox}

这些类的代码行数几乎是先前实现的两倍。但这段代码要易读得多。驱动这些类的示例也更清晰、更易扩展了：

\begin{codebox}[python]
book = Gradebook()
albert = book.get_student("Albert Einstein")
math = albert.get_subject("Math")
math.report_grade(75, 0.05)
math.report_grade(65, 0.15)
math.report_grade(70, 0.80)
gym = albert.get_subject("Gym")
gym.report_grade(100, 0.40)
gym.report_grade(85, 0.60)
print(albert.average_grade())
>>>
80.25
\end{codebox}

同样可以编写向后兼容的方法，来帮助把旧 API 风格的使用方式迁移到新的对象层次结构上。

\subsection{需要记住的要点}

\begin{itemize}
\item 避免创建这样的字典：其值是字典、长元组，或是其他内置（built-in）类型的复杂嵌套。
\item 在你需要完整类的灵活性之前，先用 \texttt{dataclasses} 这个内置（built-in）模块来构造轻量、不可变的数据容器。
\item 当你的内部状态字典变得复杂时，就把记账代码改用多个类来实现。
\end{itemize}
